\section{Introduction}
By dividing the resources among operating systems,
virtualisation is a technology that has been used
extensively over the past ten years to share the
capabilities of physical computers. IBM launched a
project in 1964 that helped define the meaning of the term
"Virtual Machines" (VM) \cite{citation_1}. The CP/CMS system was
the name for the original project, which was later changed
to Virtual Machine Facility. The Control Program
component started as an operating system back then that
allowed a computer to function as several copies of the
machine. Every copy, instance of virtual machine had its
own operating system namely the CMS component. The
concept of low level virtualisation was introduced as
described by IBM’s project and also the community effort
forwarded by independent researchers. By the middle of
the 1970s, two researchers by the name Goldberg and
Popek released multiple publications \cite{citation_2}\cite{citation_3} that provided
a comprehensive understanding of the principles and
advantages associated with the virtualisation technology.
During the next years the interest towards this new
technology decreased until VMware Inc. developed its
own product called VMware Virtual Platform for the x86-
32 architecture back in the year 1999 \cite{citation_4}. \newline \newline
A Virtual Machine Monitor, abbreviated to VMM, called
Xen \cite{citation_5} introduced the possibility of a physical computer
to boot into many operating systems within the x86-32
architecture. Its origins trace back to the academics, more
precicesly the University of Cambridge’s Computer
Laboratory. The Xen project is an open source
community project that is now actively developed by a
community with the same name. \newline \newline
The VMM implementations addressed this problem by
trapping privileged instructions with low-level
sophisticated solutions because the x86-32 architecture
was not built to facilitate full virtualisation in the way that
Popek and Goldberg outlined \cite{citation_3}. Nonetheless,
virtualisation is supported by the latest versions of the
AMD processor through its virtualisation solution
"Pacifica" \cite{citation_7}, Intel x86-32 and Itanium architectures with
extensions named VT-x for x86-32 and VT-i for Itanium
\cite{citation_6}. These features are now utilised by the majority of
virtualisation technologies, including Xen and VMware
\cite{citation_8}. \newline \newline
In a physical computer, the operating system is usually
operated in kernel mode while user applications are run in
user mode. In order for the operating system to access the
entire set of instructions, both privileged and unprivileged
of the actual processor, processors have different levels of
execution. Applications operating in the unprivileged user
mode are therefore unable to directly access the
privileged instructions. The OS's kernel mode operates in
Ring 0 of the x86 architecture's four execution modes, or
rings, whereas the user mode operates in Ring 3. Rings 1
and 2 can be utilised for further levels of access,
depending on the operating system. \newline \newline
The execution of whole operating systems can currently
be supported by a number of virtualisation approaches. In
contrast to Process VMs, such as those provided by the
Java Virtual Machine(JVM), VMs for the execution of
guest operating systems are also known as System VMs
\cite{citation_10}. Going forward with the paper, I’ll provide an
overview of the virtualisation strategies from an operating
systems perspective in the following.